\section{Approach}
\label{sec:approach}
Our approach rests on one premise: although flakiness manifests only at run time, it leaves a stable lexical trace in the test code. Figure~\ref{fig:metodologia} summarizes the six-step pipeline.

\begin{figure*}[!htbp]
\centering
{\scriptsize
\setlength{\fboxsep}{3pt}%
\newcommand{\etapa}[2]{\fbox{\parbox[c][13.5ex][c]{\dimexpr(\textwidth-55pt)/6-2\fboxsep-2\fboxrule\relax}{\centering\hyphenpenalty=10000\textbf{#1}\\[2pt]#2}}}%
\etapa{1. Subjects}{17 open-source\\C\# projects;\\19,211 test cases}%
\hfill$\rightarrow$\hfill
\etapa{2. Reruns}{100$\times$ + 20$\times$\\isolated;\\TRX reports}%
\hfill$\rightarrow$\hfill
\etapa{3. Labels}{flaky, broken,\\stable, skipped;\\12,665 methods,\\30 flaky}%
\hfill$\rightarrow$\hfill
\etapa{4. Processing}{Roslyn: tokens\\and structure}%
\hfill$\rightarrow$\hfill
\etapa{5. Models}{TF-IDF;\\8 classifiers;\\1:1; CV and LOPO}%
\hfill$\rightarrow$\hfill
\etapa{6. Vocabulary}{information\\gain of\\tokens}%
}
\caption{End-to-end pipeline for building the vocabulary of flaky tests in C\#.}
\label{fig:metodologia}
\end{figure*}

\myparagraph{Step 1: Project selection} We started from open-source C\# libraries hosted on GitHub with extensive test suites and then filtered them by an objective criterion: the project had to build and run its entire suite in our environment. \Den{Do jeito que está, a seleção não é reproduzível. Faltam: como os candidatos foram encontrados (busca no GitHub? lista curada?), quais critérios (estrelas, atividade, número de testes), quantos candidatos foram tentados e quantos falharam no filtro, e em qual commit (SHA) cada projeto foi analisado. O filtro ``builda e roda tudo'' gera viés de sobrevivência para suítes fáceis de executar.} The result is 17 projects and 19,211 test cases, covering messaging, actor systems, scheduling, RPC, logging, resilience, and architecture testing. Two candidate projects (StackExchange.Redis and SSH.NET) were excluded because their suites depend on external servers whose runs were disrupted by the environment.

\myparagraph{Step 2: Reruns} Each suite was built once and run 100 times in a row, without rebuilding between runs. \Den{A ordem dos testes é a mesma em todas as execuções? Se for, a flakiness dependente de ordem (OD) não é exposta e o estudo cobre essencialmente NOD; dizer isso explicitamente, como Pinto et al. Informar também a configuração de paralelismo (o xUnit paraleliza coleções por padrão, o NUnit não), os timeouts e o que acontece quando o runner trava.} In each run, the .NET test runner (\texttt{dotnet test}) produced a TRX (Test Results XML) report, an XML file recording each executed test. We built with \texttt{dotnet build} and ran with \texttt{dotnet test --no-build} and the TRX logger; xUnit v3 projects on Microsoft.Testing.Platform were run through the test executable itself, with \texttt{--report-trx}. These reports were the basis for identifying tests whose outcomes changed across runs. Some rounds suffered environmental disruptions, such as leftover containers from previous runs and a degraded Redis cluster. They were discarded and repeated in isolation. Since flakiness is rare (fewer than 1\% of tests), a few unstable tests are enough to change the results, and any environmental noise could be mistaken for flakiness. We therefore performed a second round of 20 runs under controlled isolation (5 for Testcontainers, whose suite requires containers) to check the tests flagged in the first and to look for flaky tests that the first round missed. Since a single pass/fail flip on the same version already shows that a test is flaky, the second round can add flaky tests but cannot refute those found in the first; Section~\ref{sec:rq12} compares the two rounds. Everything ran on a laptop with an Intel Core i7-1165G7, 16 GB of RAM, and Windows 11. \Den{Sugestão (requer execução): uma rodada adicional em CI (GitHub Actions/Linux) ou com stress de CPU (estilo Shaker) aumentaria o número de positivos e responderia à crítica sobre máquina única. A limitação já foi registrada nas ameaças (Seção~\ref{sec:threats}).} The 100-run rounds took place in July 2026, before a .NET SDK update whose previous version we did not record; the 20-run rounds used the SDK resolved by each repository's \texttt{global.json} (Section~\ref{sec:threats} discusses the consequences). \Den{Definir o que é ``controlled isolation'' (máquina ociosa? um repositório por vez? rede desligada?). Se possível, mostrar que o impacto do SDK não registrado é pequeno (o \texttt{global.json} de vários repositórios fixa o SDK). A frase sobre a limitação foi movida para as ameaças.} The SDKs installed side by side are 8.0.425, 9.0.318, 10.0.112, 10.0.303, and 10.0.401; the 20-run round used 10.0.401 in 11 repositories, 10.0.303 in Nerdbank.Streams, Polly, and StreamJsonRpc, 10.0.112 in Akka.NET, and 8.0.425 in the RabbitMQ client.

\myparagraph{Step 3: Labeling} Each test received a label according to its outcomes across runs: flaky (passes and fails), broken (always fails), stable (never fails), or skipped (skipped in every run). \Den{Definir os casos de borda: teste que passa em algumas execuções e é pulado em outras, timeouts, crash do processo de teste (todos os testes daquela execução viram ``failed''?), testes ausentes em algumas TRX. Cada um muda o rótulo.} Broken and skipped tests were excluded from modeling, which compares only flaky against stable tests. \Den{Qual rótulo é usado na modelagem quando as rodadas discordam? Pelo texto, os 6 testes flaky só na rodada de 20 entram como \emph{stable} no conjunto de 100, o que é ruído de rótulo conhecido. Sugiro usar o rótulo união (36 flaky) como principal ou, no mínimo, como análise de sensibilidade.} The unit of analysis is the declaring test method, not the test case reported by the runner. Parameterized tests and methods inherited from base classes generate several cases from a single method. Without this grouping, the same code would be counted several times, and nearly identical tests would fall into different cross-validation folds. The choice of unit changes the count: the 30 flaky methods correspond to 50 flaky test cases, because some methods generate several cases (through parameterization or inheritance). Since test identity in TRX files is fragile, we identified each test by its fully qualified class plus its display name and handled the pitfalls in Table~\ref{tab:armadilhas}: we normalized timestamps, thread identifiers, and run dates in the names, ran both rounds on the same target framework, and discarded contaminated rounds.

\Den{Tab.~\ref{tab:armadilhas}: os números de herdados não batem com a Tab.~\ref{tab:lopo-repo}. Nerdbank.Streams: 934 métodos com 363 herdados dariam 571 métodos declarantes, mas a Tab.~\ref{tab:lopo-repo} lista 595 testes. StreamJsonRpc: 2.495 $-$ 2.199 = 296, mas a Tab.~\ref{tab:lopo-repo} lista 763. Explicar como cada número é contado ou corrigir.}
\begin{table}[!htbp]
\caption{Methodological pitfalls of rerun studies in .NET.}
\label{tab:armadilhas}
\centering
\footnotesize
\renewcommand{\arraystretch}{1.15}
\begin{tabular}{@{}p{0.30\columnwidth}p{\dimexpr0.70\columnwidth-2\tabcolsep\relax}@{}}
\toprule
\raggedright \textbf{Pitfall} & \raggedright \textbf{Example and effect} \tabularnewline
\midrule
\raggedright Nondeterministic names & \raggedright Akka.NET \emph{theories} embed timestamps and \emph{thread} ids: 3,308 distinct names for 1,328 tests \tabularnewline
\raggedright Dates in names & \raggedright NLog \emph{theories} embed the current date: 2,876 vs.\ 2,870 tests across rounds \tabularnewline
\raggedright Target framework & \raggedright Nerdbank.Streams: 1,142 tests on \texttt{net472} vs.\ 877 on \texttt{net8.0} \tabularnewline
\raggedright Name collisions & \raggedright In some projects, TRX stores only the method name: tests from different classes get mixed \tabularnewline
\raggedright Contaminated environment & \raggedright Leftover containers and a degraded \emph{cluster} produce spurious flaky tests \tabularnewline
\raggedright Inherited tests & \raggedright 363 of 934 methods (Nerdbank.Streams) and 2,199 of 2,495 (StreamJsonRpc) are inherited \tabularnewline
\bottomrule
\end{tabular}
\end{table}

\myparagraph{Step 4: Test processing} Using Roslyn, the C\# compiler, we extracted each test method from the source code, identifying it by its attributes (\texttt{[Fact]}, \texttt{[Test]}, \texttt{[TestMethod]}, and equivalents). From each method, we recorded the tokens (identifiers and keywords), discarding reserved words and common language modifiers (such as \texttt{public}, \texttt{static}, \texttt{var}, \texttt{if}, and \texttt{return}), which appear in almost every test and do not distinguish them. \Den{Detalhar a tokenização: os identificadores são quebrados em camelCase/PascalCase (\texttt{StopAsync} $\to$ \texttt{stop}, \texttt{async})? Pelos tokens reportados (\texttt{checkgcpressureasync}, \texttt{cancellationtokensource}) parece que não, o que cria tokens únicos por projeto e reforça a especificidade de projeto. Testar as duas versões. Dizer também se strings, comentários e nomes de métodos/atributos entram, e se há lowercasing/stemming.} We kept \texttt{async} and \texttt{await}, which indicate asynchronous code. We also recorded size and structural features: number of lines, number of \texttt{await} expressions and assertions, and whether the method is \texttt{async}. Two of these features count keywords, so the structure-only baseline (Step~5) is not entirely free of lexical information. \Den{Deixar claro se os modelos principais usam só TF-IDF ou TF-IDF + estrutura. Se possível, avaliar um modelo combinado e um baseline estrutural puro (só tamanho/asserções), o que requer execução.} Since many tests are inherited, we resolved each executed test to its declaring method, which is flaky if any inherited instance is flaky. \Den{Atenção: quando um teste herdado só é flaky numa classe derivada (por exemplo, com um serializer específico no StreamJsonRpc), a causa provavelmente está no \emph{setup} da derivada, não no corpo do método declarante, e o vocabulário do corpo não tem como explicá-la. Reportar quantos dos 30 flaky estão nessa situação.}

\myparagraph{Step 5: Model construction} The tokens form the test's vocabulary and are vectorized with TF-IDF (term frequency--inverse document frequency), a technique that turns the text of each test into a numeric vector. Each term receives a weight that combines how often it appears in that test (TF) and the inverse of the proportion of tests that contain it (IDF). Thus, the weight is high when the term appears often in the test and is rare in the set (such as \texttt{weakreference}) and low when it is present in almost every test (such as \texttt{assert}). We used single terms (unigrams, such as \texttt{await}) and pairs of adjacent terms (bigrams, such as \texttt{await task}), with sublinear term frequency, which reduces the weight of heavily repeated terms. We kept only terms present in at least two tests and in at most 90\% of them. The vectorizer is fitted only on the training tests of each fold, which prevents information leakage from the evaluation tests. \Den{Com 30 positivos e undersampling 1:1, cada fold treina com $\approx$24 flaky + 24 estáveis e testa com 6 + 6: milhares de features TF-IDF para 48 exemplos, e um MCC estimado sobre 12 testes. É a principal fragilidade estatística do artigo. Alternativas: (i) usar todos os estáveis com \texttt{class\_weight='balanced'} ou fazer undersampling só no treino, avaliando em folds com a taxa natural; (ii) usar \texttt{StratifiedGroupKFold} agrupando por classe/arquivo, para que testes quase idênticos (ex.: \texttt{JsonRpcPerf\_Pipe} e \texttt{JsonRpcPerf\_Channel}) não caiam em folds diferentes.} We trained eight classifiers: Random Forest and Extra Trees (200 trees each), Gradient Boosting, Decision Tree, Naive Bayes, linear SVM, KNN, and logistic regression. \Den{Os estáveis são sorteados de todos os 17 projetos, inclusive dos 8 sem nenhum flaky. Assim o modelo pode aprender ``de qual projeto vem o teste'' (convenções, framework, nomes) em vez de flakiness. Exemplo: \texttt{areequal} (NUnit) aparece em 20\% dos estáveis e 0\% dos flaky (Seção~\ref{sec:rq3}). Sugestões: (i) um baseline que vê só o nome do repositório (no paper de Go ele empata com o vocabulário); (ii) sortear os estáveis dos mesmos repositórios dos flaky (amostragem pareada). Sem isso, um revisor que conheça o paper de Go ou o de Gültekin et al. vai apontar \emph{project leakage}.} The first three are ensemble models, which combine many decision trees. Apart from the settings listed in Section~\ref{sec:setup}, we used default hyperparameters and did not tune them, since tuning on 30 positives would add a further source of overfitting. To find out what the vocabulary actually adds, we compared all of them against four baselines: a majority predictor, a random predictor, a threshold rule based on the 10 tokens with the highest information gain, \Den{Onde o IG da regra é calculado? Se for no conjunto inteiro (como no Step 6, com 30 subamostras de todos os dados), há vazamento: a regra viu os testes de avaliação. Ela precisa ser recalculada no treino de cada fold/LOPO. Explicar também o limiar (basta qualquer um dos 10 tokens estar presente?).} and a Random Forest trained only on the size and structural features, without any tokens.

\myparagraph{Step 6: Vocabulary extraction} Finally, we ranked the tokens by information gain with respect to the label, which measures how much knowing whether the token is present reduces uncertainty about the test's class, in nats and averaged over 30 subsampling draws, together with its frequency in each class, the number of repositories in which it occurs in flaky tests, and the lift between classes, that is, how many times more frequent the token is among flaky tests than among stable ones. This ranking makes the approach explanatory rather than opaque: it allows the predictive signal to be linked to concrete constructs and to the known root causes of flakiness. \Den{Sugestão: complementar com o IG leave-one-project-out do paper de Go (quanto do IG de cada token sobra quando o projeto dominante é removido). É um diagnóstico barato e responde diretamente à observação de que o ranking bruto é dominado por StreamJsonRpc (Seção~\ref{sec:rq3}).}